Kiesling et al.\ already combine security-control models, simulation,
multi-objective optimization, and portfolio selection \cite{kiesling2016}.
Our hospital-specific contribution is the explicit evidence-to-mechanism
framework, clinical-service dependency outcomes, exhaustive paired finite
comparison, and separation of endpoint retention from other uncertainties.
Hospital process simulation with local observations and expert input
\cite{willing2025} addresses a different validation level; our synthetic
network does not inherit that empirical grounding.

Necessary and possible efficiency under interval uncertainty are established
optimization concepts \cite{hladik2017}. Equations~\ref{eq:guaranteed}--\ref{eq:possible}
give a direct finite-box specialization, proof, and witness-producing
implementation. The added evidence is its endpoint application and
transparent evaluation on the published RE suite \cite{tanabe2020}.
The benchmark perturbations assess this implementation, not the continuous
engineering problems' global Pareto fronts.

STRESS reporting guidance \cite{monks2019}, common random numbers
\cite{buffalo2026}, and the distinction between prior plans and retrospective
analysis \cite{nosek2018} inform the experimental design. Multiverse and
specification-curve work \cite{steegen2016,simonsohn2020} motivate examining
analytical choices; our enumeration is descriptive and does not implement
specification-curve joint inference.

Statistical consistency checks \cite{nuijten2016}, metamorphic simulation
tests \cite{raunak2021}, and the recent ABE-Ralph experimental-fidelity
preprint \cite{yu2026} already address relations beyond successful execution.
Our small deterministic registry contributes case-specific scope, unit,
and quantifier checks. Its role is to make interpretations reviewable beside
the model and decision diagnostics, rather than to claim general automated
scientific judgment.
